% =====================================================================
%  HÌNH 0 (FLOW) — CƠ CHẾ: vì sao cùng một lượng byte lại cho số vòng khác nhau
%
%  Đây là hình KHÁI NIỆM của bài. Nó phải trả lời đúng một câu: hai giao thức
%  mang CÙNG khối lượng khoá hậu lượng tử, nhưng một bên trả bằng SỐ VÒNG còn
%  bên kia trả bằng ĐỘ DÀY của flight.
%
%  ⚠ MỌI CON SỐ ở đây đến từ figures/out/numbers.tex, do make_figures.py sinh
%  ra từ analysis/model.py. KHÔNG gõ số vào tệp này: gõ tay là tạo chỗ ở thứ
%  hai cho cùng một con số, và chạy lại mô hình sẽ không cập nhật hình.
%
%  Dịch: pdflatex fig0-flow.tex   (chạy từ thư mục figures/)
% =====================================================================
\documentclass[border=3pt]{standalone}
\usepackage{tikz}
\usepackage{amsmath}
\input{../../../transaction-figure-kit/flow/txstyle.tex}
\usetikzlibrary{decorations.pathreplacing}
%% SINH TU figures/make_figures.py -- KHONG sua tay.
\newcommand{\numFramePayload}{102}
\newcommand{\numCoapBlock}{64}
\newcommand{\numDtlsRT}{2}
\newcommand{\numMsgOneClassic}{37}
\newcommand{\numMsgOnePQ}{1192}
\newcommand{\numMsgTwoPQ}{4415}
\newcommand{\numExchOnePQ}{19}
\newcommand{\numExchTwoPQ}{69}
\newcommand{\numExchTotalPQ}{140}
\newcommand{\numKemFiveOneTwo}{800}
\newcommand{\numMOneMatch}{7}
\newcommand{\numMOneTotal}{7}
\newcommand{\numMOneImpl}{aiocoap 0.4.17}
\newcommand{\numGnutlsOk}{17}
\newcommand{\numGnutlsTotal}{21}
\newcommand{\numGnutlsAtFrame}{1/3}
\newcommand{\numGnutlsFragOk}{23}
\newcommand{\numGnutlsFragBad}{28}
\newcommand{\numMbedtlsOk}{14}
\newcommand{\numMbedtlsTotal}{14}
\newcommand{\numMbedtlsAtFrame}{2/2}
\newcommand{\numOpensslOk}{9}
\newcommand{\numOpensslTotal}{21}
\newcommand{\numOpensslAtFrame}{0/3}
\newcommand{\numOpensslMtuBad}{250}
\newcommand{\numOpensslMtuOk}{300}
\newcommand{\numKemFrameRatio}{7.8}
\newcommand{\numDecompError}{4.2}
\newcommand{\numBestBlock}{1024}
\newcommand{\numBestExch}{11}
\newcommand{\numBestRatio}{5.5}
\newcommand{\numQBlockMaxPayloads}{10}
\newcommand{\numQBlockExch}{15}
\newcommand{\numQBlockGain}{9.3}
\newcommand{\numFramePayloadSecured}{81}
\newcommand{\numMacOverhead}{25}
\newcommand{\numFramePayloadFrame}{127}
\newcommand{\numLinkSec}{21}
\newcommand{\numFragHdr}{5}
\newcommand{\numExchLo}{135}
\newcommand{\numExchHi}{147}
\newcommand{\numDirDefault}{280}
\newcommand{\numDirBest}{22}
\newcommand{\numDirQBlock}{30}
\newcommand{\numRatioDefault}{47}
\newcommand{\numRatioBest}{4}
\newcommand{\numRatioQBlock}{5}
\newcommand{\numFramesDefault}{280}
\newcommand{\numRatioFramesDefault}{2.43}
\newcommand{\numRatioFrames}{1.15}
\newcommand{\numBestFrames}{132}
\newcommand{\numGnutlsCapKB}{2.3}
\newcommand{\numMbedtlsMaxFrag}{32}
\newcommand{\numOpensslFewFragFail}{5}
\newcommand{\numOpensslManyFragOk}{19}
\newcommand{\numDtlsMeasTurns}{6}
\newcommand{\numDtlsMaxFrag}{94}
\newcommand{\numDtlsMinFrag}{32}
\newcommand{\numDtlsDatagramLo}{55}
\newcommand{\numDtlsDatagramHi}{115}
\newcommand{\numDtlsDatagramGrowth}{2.09}
\newcommand{\numDtlsImpl}{mbedtls 3.6.2}
\newcommand{\numDtlsMeasVersion}{1.2}
\newcommand{\numChainKeyBits}{8192}
\newcommand{\numSpanOk}{22.4}
\newcommand{\numSpanBad}{22.95}
\newcommand{\numSpanCells}{30}
\newcommand{\numSpanCerts}{3}
\newcommand{\numSpanMtus}{10}
\newcommand{\numSpanTrials}{150}
\newcommand{\numSpanCapKB}{2.3}
\newcommand{\numAmbiguousFrag}{23}
\newcommand{\numLongAllowance}{240}
\newcommand{\numLongRuns}{18}
\newcommand{\numLongSelfExit}{18}
\newcommand{\numLongKilled}{0}
\newcommand{\numLongSeconds}{46}
\newcommand{\numRelayCtrlCells}{30}


\tikzset{
  actor/.style={blk, minimum width=19mm, minimum height=7mm, font=\scriptsize},
  life/.style={draw=inkgray!55, line width=0.5pt, dash pattern=on 2pt off 2pt},
  msg/.style={->, draw=inkgray, line width=0.5pt},
  cmsg/.style={->, draw=accent, line width=0.7pt},
  ackm/.style={->, draw=inkgray!60, line width=0.45pt, dash pattern=on 2pt off 1.5pt},
  mlab/.style={font=\tiny, text=inkgray, fill=white, inner sep=1pt},
  rt/.style={font=\tiny\itshape, text=accent},
  panel/.style={draw=inkgray!35, line width=0.4pt, rounded corners=1pt,
                inner sep=3.5mm},
  ptitle/.style={font=\scriptsize\bfseries, text=inkgray},
}

\begin{document}
\begin{tikzpicture}[x=1cm, y=1cm]

% =====================================================================
% (a) EDHOC trên CoAP: block-wise LOCK-STEP — mỗi khối MỘT vòng
% =====================================================================
\begin{scope}[shift={(0,0)}]
  \node[actor] (I) at (0,0)   {Initiator};
  \node[actor] (R) at (4.2,0) {Responder};
  \foreach \a in {I,R} \draw[life] (\a.south) -- ++(0,-6.2);

  % message_1, cắt khối
  \node[rt, anchor=west] at (-1.85,-0.55) {message\_1};
  \draw[cmsg] (I|-0,-0.9) -- node[mlab]{Block2 \#0 (\numCoapBlock\,B)} (R|-0,-0.9);
  \draw[ackm] (R|-0,-1.35) -- node[mlab]{ACK} (I|-0,-1.35);
  \draw[cmsg] (I|-0,-1.8) -- node[mlab]{Block2 \#1} (R|-0,-1.8);
  \draw[ackm] (R|-0,-2.25) -- node[mlab]{ACK} (I|-0,-2.25);
  \node[font=\scriptsize, text=inkgray] at (2.1,-2.72) {$\vdots$};
  \draw[cmsg] (I|-0,-3.15) -- node[mlab]{Block2 \#\the\numexpr\numExchOnePQ-1\relax}
              (R|-0,-3.15);
  \draw[ackm] (R|-0,-3.6) -- node[mlab]{ACK} (I|-0,-3.6);

  % ngoặc: số vòng của message_1
  \draw[decorate, decoration={brace, amplitude=3pt, mirror}, inkgray!70,
        line width=0.4pt] (4.75,-0.9) -- (4.75,-3.6)
        node[midway, right=3pt, rt, align=left]
        {\numExchOnePQ{} round trips\\ \tiny(\numMsgOnePQ\,B $\div$ \numCoapBlock\,B)};

  % message_2 và _3, rút gọn
  \node[rt, anchor=west] at (-1.85,-4.15) {message\_2};
  \draw[cmsg] (R|-0,-4.45) -- node[mlab]{\numExchTwoPQ{} blocks, lock-step} (I|-0,-4.45);
  \node[rt, anchor=west] at (-1.85,-5.15) {message\_3};
  \draw[cmsg] (I|-0,-5.45) -- node[mlab]{\dots} (R|-0,-5.45);

  \node[ptitle, align=center] at (2.1,0.85)
    {(a) EDHOC over CoAP\\[-1pt] \tiny block-wise, LOCK-STEP [RFC 7959]};
  \node[font=\scriptsize, text=accent, align=center] at (2.1,-6.75)
    {TOTAL \numExchTotalPQ{} round trips};
\end{scope}

% =====================================================================
% (b) DTLS 1.3: theo FLIGHT — nhiều datagram, vẫn hai vòng
% =====================================================================
\begin{scope}[shift={(9.4,0)}]
  \node[actor] (C) at (0,0)   {Client};
  \node[actor] (S) at (4.2,0) {Server};
  \foreach \a in {C,S} \draw[life] (\a.south) -- ++(0,-6.2);

  % flight 1: nhiều bản ghi bắn liên tiếp, KHÔNG chờ
  \node[rt, anchor=west] at (-1.75,-0.55) {flight 1};
  \foreach \y in {0.9, 1.12, 1.34, 1.56} {
    \draw[cmsg] (C|-0,-\y) -- (S|-0,-\y);
  }
  \node[mlab, anchor=west] at (0.35,-1.78)
    {several MTU-sized records, sent BACK-TO-BACK};
  \draw[decorate, decoration={brace, amplitude=3pt, mirror}, inkgray!70,
        line width=0.4pt] (4.75,-0.9) -- (4.75,-1.56)
        node[midway, right=3pt, rt, align=left] {1 round trip};

  % flight 2
  \node[rt, anchor=west] at (-1.75,-2.75) {flight 2};
  \foreach \y in {3.1, 3.32, 3.54, 3.76, 3.98} {
    \draw[cmsg] (S|-0,-\y) -- (C|-0,-\y);
  }
  \draw[decorate, decoration={brace, amplitude=3pt, mirror}, inkgray!70,
        line width=0.4pt] (4.75,-3.1) -- (4.75,-3.98)
        node[midway, right=3pt, rt, align=left] {1 round trip};

  % flight 3
  \node[rt, anchor=west] at (-1.75,-4.6) {flight 3};
  \foreach \y in {4.9, 5.12} {
    \draw[cmsg] (C|-0,-\y) -- (S|-0,-\y);
  }

  \node[ptitle, align=center] at (2.1,0.85)
    {(b) DTLS 1.3\\[-1pt] \tiny fragments at the handshake layer, by FLIGHT [RFC 9147]};
  \node[font=\scriptsize, text=accent, align=center] at (2.1,-6.75)
    {TOTAL \numDtlsRT{} round trips, INDEPENDENT of size};
\end{scope}

% =====================================================================
% Câu chốt: cùng lượng byte, hai cách trả giá
% =====================================================================
\node[panel, fill=accent!4, draw=accent!45, text width=15.5cm, align=center,
      font=\scriptsize] at (6.7,-7.9)
  {Both carry the \emph{same} volume of post-quantum key material. Side (a) pays
   in \textbf{round trips}, because each CoAP block is its own request--response
   pair; side (b) pays in \textbf{flight thickness}, and its round-trip count does
   not move. Classically \emph{both} take \numDtlsRT{} round trips: an EDHOC
   message fits inside one \numFramePayload\,B frame, so block-wise never engages.
   Post-quantum credentials break exactly that condition --- the ML-KEM-512
   encapsulation key alone is \numKemFiveOneTwo\,B.};

\end{tikzpicture}
\end{document}
